\begin{table}[htbp]
\centering
\caption{Standardized Effect Sizes}
\label{tab:sde}
\begin{adjustbox}{max width=\textwidth}
\begin{threeparttable}
\begin{tabular}{l cccccc}
\toprule
Outcome & $\hat{\beta}$ & SE & SD($Y$) & SDE & SE(SDE) & Classification \\
\midrule
\multicolumn{7}{l}{\textit{Panel A: Pooled}} \\
CC Total Enrollment & -0.0108 & 0.0042 & 1.266 & -0.0085 & 0.0033 & Small negative \\
CC Black Enrollment & -0.0131 & 0.0054 & 2.055 & -0.0064 & 0.0026 & Small negative \\
CC Hispanic Enrollment & -0.0227 & 0.0065 & 2.094 & -0.0108 & 0.0031 & Small negative \\
CC White Enrollment & -0.0150 & 0.0047 & 1.509 & -0.0100 & 0.0031 & Small negative \\
\midrule
\multicolumn{7}{l}{\textit{Panel B: Heterogeneous (by displacement intensity)}} \\
CC Total (High Exposure) & -0.0135 & 0.0056 & 1.266 & -0.0107 & 0.0044 & Small negative \\
CC Total (Low Exposure) & -0.0027 & 0.0071 & 1.266 & -0.0021 & 0.0056 & Null \\
\bottomrule
\end{tabular}
\begin{tablenotes}[flushleft]
\footnotesize
\item \textit{Notes:} \textbf{Country:} United States. \textbf{Research question:} Does the mass closure of for-profit colleges following the 2016 ACICS accreditation revocation increase enrollment at nearby community colleges, and does absorption differ by race? \textbf{Policy mechanism:} The Department of Education revoked recognition of ACICS, the largest for-profit accreditor, in September 2016, causing approximately 247 institutions enrolling 600,000 students to lose Title IV federal financial aid eligibility and close within 18 months, displacing students who must find alternative educational institutions or exit higher education entirely. \textbf{Outcome definition:} Log undergraduate enrollment at public 2-year institutions (community colleges) aggregated to the county-year level, from IPEDS enrollment tables. \textbf{Treatment:} Continuous: log(1 + total 2015 enrollment at for-profit institutions in the county that closed by 2018). \textbf{Data:} IPEDS institution-year panel, 2010--2022, aggregated to county-year level; approximately 1,200 counties with community colleges over 13 years. \textbf{Method:} Two-way fixed effects DiD with county and year fixed effects; standard errors clustered at the state level. \textbf{Sample:} Counties with at least one public 2-year institution; treatment counties contain at least one for-profit institution that closed between 2015 and 2018. SDE $= \hat{\beta} / \text{SD}(Y)$ where SD($Y$) is the pre-treatment standard deviation. Classification refers to magnitude, not statistical significance: Large ($|$SDE$| > 0.15$), Moderate ($0.05$--$0.15$), Small ($0.005$--$0.05$), Null ($< 0.005$).
\end{tablenotes}
\end{threeparttable}
\end{adjustbox}
\end{table}

